
Data curation recipes for pre-training language models are typically developed at a single model scale and applied uniformly across model sizes. This paper investigates whether the optimal perplexity-based filtering threshold for pre-training data is a function of model capacity. We conduct a controlled factorial experiment varying GPT-2-scored perplexity filtering threshold $\tau \in$ {20, 35, 50} and MinHash deduplication aggressiveness (Jaccard threshold J $\in$ {0.7, 0.9}) across two model scales (approximately 7.9M and 18.1M parameters) on FineWeb (sample-10BT). Across all 24 experimental runs (3 PPL $\times$ 2 dedup $\times$ 2 scales $\times$ 2 seeds), the 7.9M-parameter model achieves its highest HellaSwag 0-shot normalized accuracy under strict filtering ($\tau =20$; 0.2556), while the 18.1M-parameter model achieves its highest accuracy under permissive filtering ($\tau =50$; 0.2548) --- with strict filtering being the worst configuration for the larger model (0.2524). The effect magnitude is small ($\Delta acc$\_norm $\approx$ 0.003) but directionally consistent across all runs. Both models exceed the random baseline of 0.25 in all 24 conditions. These results provide preliminary evidence that the optimal pre-training curation configuration is scale-dependent, though the experiments operate far below the originally targeted scales (70M/160M parameters) and training duration (50B tokens), and formal statistical significance was not established at this PoC scale. Deduplication effects across scales and learning curve dynamics remain uncharacterized due to compute and storage constraints.
